% ============================================================
% abstract.tex  [ENGLISH — RCR submission version]
% Paper: Municipal Public Cleaning Expenditure and Solid Waste
%        Management: Evidence from Peru, 2015–2019
% Author: Dante Barreto Gamarra (UNAS)
% Note: 148 words — complies with INV-5 (150-word limit)
% ============================================================

\begin{abstract}
\noindent \singlespacing
This study estimates the effect of municipal public cleaning expenditure on solid waste management across Peruvian districts from 2015 to 2019. The main identification threat is expenditure endogeneity: municipalities with stronger institutional capacity simultaneously execute higher budgets and report better outcomes. To correct this bias, we apply the Correlated Random Effects with Control Function (CRE+CF) estimator proposed by \citet{Wooldridge2015}, combining the \citet{Mundlak1978} device with first-stage residuals as an additional regressor. The analysis covers three management dimensions using a panel of 9,322 observations, 1,874 municipalities, and 196 provincial clusters. A 10\% increase in accrued expenditure raises daily waste collection by 8.79\% (elasticity $= +0.879$; $p < 0.001$), increases the municipal planning index by 0.030 points on a five-instrument scale ($\hat{\beta} = +0.300$; $p < 0.001$), and raises the probability of formal sanitary landfill disposal by 2.2 percentage points over an 18\% baseline (AME $= +0.022$; $p < 0.001$). All three primary hypotheses survive Bonferroni-Holm correction. Results identify municipal finance as the binding constraint for environmental formalization of final disposal, and support reforms to Peru's FONCOMUN transfer formula targeting municipalities with the largest solid waste management deficits.
\end{abstract}

\vspace{1em}
\noindent \textbf{JEL Codes:} H72, Q53, R51, O18, H77

\vspace{0.5em}
\noindent \textbf{Keywords:} municipal expenditure, solid waste management, sanitary landfill, CRE+CF, Mundlak device, Peru, fiscal decentralization, FONCOMUN
